% Bagian Tinjauan Pustaka: Static Cycle Counting
% Untuk dimasukkan ke dalam Bab 2

\subsection{Static Cycle Counting untuk Estimasi Waktu Inferensi}
\label{subsec:static-cycle-counting}

Evaluasi kinerja temporal model \textit{machine learning} pada mikrokontroler memerlukan metode pengukuran yang akurat dan dapat direproduksi. Pada penelitian ini, digunakan metode \textit{static cycle counting} untuk mengestimasi waktu inferensi model pada platform MSP430F5529. Metode ini dipilih karena memberikan hasil yang deterministik dan tidak memerlukan perangkat keras target secara langsung.

\subsubsection{Definisi dan Prinsip Dasar}

\textit{Static cycle counting} adalah teknik analisis statis yang menghitung jumlah siklus CPU yang diperlukan untuk mengeksekusi suatu program dengan menganalisis kode assembly hasil kompilasi tanpa menjalankan program tersebut secara aktual \cite{wilhelm2008wcet}. Berbeda dengan metode pengukuran dinamis yang mengukur waktu eksekusi secara langsung pada perangkat keras, metode statis menganalisis setiap instruksi dalam kode mesin dan menjumlahkan biaya siklus berdasarkan spesifikasi arsitektur prosesor.

Prinsip dasar metode ini adalah bahwa setiap instruksi pada arsitektur prosesor tertentu memiliki jumlah siklus eksekusi yang telah ditentukan oleh pabrikan. Untuk arsitektur MSP430, informasi ini tersedia dalam dokumen \textit{MSP430x5xx and MSP430x6xx Family User's Guide} \cite{ti2018msp430}. Dengan menganalisis kode assembly menggunakan \textit{disassembler} seperti \texttt{msp430-elf-objdump}, setiap instruksi dapat diidentifikasi dan biaya siklusnya dapat dihitung.

\subsubsection{Siklus Instruksi MSP430}

Arsitektur MSP430 menggunakan set instruksi RISC (\textit{Reduced Instruction Set Computer}) 16-bit dengan jumlah siklus yang bervariasi tergantung pada jenis instruksi dan mode pengalamatan yang digunakan. Tabel \ref{tab:msp430-cycles} menunjukkan jumlah siklus untuk instruksi-instruksi utama yang relevan dengan algoritma inferensi pohon keputusan.

\begin{table}[htbp]
\centering
\caption{Jumlah Siklus Instruksi MSP430 yang Relevan}
\label{tab:msp430-cycles}
\begin{tabular}{|l|c|l|}
\hline
\textbf{Instruksi} & \textbf{Siklus} & \textbf{Keterangan} \\
\hline
\texttt{MOV Rn, Rm} & 1 & Salin register ke register \\
\texttt{MOV @Rn, Rm} & 2 & Muat dari memori (indirect) \\
\texttt{MOV X(Rn), Rm} & 3 & Muat dari memori (indexed) \\
\texttt{ADD Rn, Rm} & 1 & Penjumlahan register \\
\texttt{CMP Rn, Rm} & 1 & Perbandingan register \\
\texttt{CMP @Rn, Rm} & 2 & Perbandingan dengan memori \\
\texttt{JNE/JEQ/JMP} & 2 & Lompatan kondisional/tak kondisional \\
\texttt{CALL} & 4 & Pemanggilan fungsi \\
\texttt{RET} & 3 & Kembali dari fungsi \\
\texttt{PUSH} & 3 & Simpan ke stack \\
\texttt{POP} & 2 & Ambil dari stack \\
\hline
\end{tabular}
\end{table}

Mode pengalamatan mempengaruhi jumlah siklus tambahan. Pengalamatan \textit{indexed} \texttt{X(Rn)} menambah 2 siklus, pengalamatan \textit{indirect} \texttt{@Rn} menambah 1 siklus, dan pengalamatan \textit{immediate} \texttt{\#N} menambah 1 siklus untuk operand sumber.

\subsubsection{Metodologi Penghitungan}

Proses \textit{static cycle counting} pada penelitian ini dilakukan melalui tahapan berikut:

\begin{enumerate}
    \item \textbf{Kompilasi Kode C}: Kode inferensi model dikompilasi menggunakan \texttt{msp430-elf-gcc} dengan optimisasi level \texttt{-O2} untuk menghasilkan file ELF (\textit{Executable and Linkable Format}).
    
    \item \textbf{Disassembly}: File ELF di-\textit{disassemble} menggunakan \texttt{msp430-elf-objdump -d} untuk mendapatkan kode assembly beserta alamat dan opcode setiap instruksi.
    
    \item \textbf{Identifikasi Fungsi}: Batas-batas fungsi kritis (\texttt{tree\_predict}, \texttt{ensemble\_predict}, \texttt{quantize\_features}) diidentifikasi dari tabel simbol menggunakan \texttt{msp430-elf-nm}.
    
    \item \textbf{Penghitungan Siklus}: Setiap instruksi dalam fungsi-fungsi tersebut dianalisis dan biaya siklusnya dijumlahkan berdasarkan tabel referensi instruksi MSP430.
    
    \item \textbf{Estimasi Jalur Eksekusi}: Untuk struktur kontrol seperti loop dan percabangan, digunakan estimasi jalur rata-rata berdasarkan karakteristik algoritma (misalnya, kedalaman pohon rata-rata).
\end{enumerate}

\subsubsection{Model Analitis untuk Inferensi Pohon Keputusan}

Berdasarkan analisis struktur algoritma inferensi ensemble pohon keputusan, total siklus dapat dimodelkan sebagai:

\begin{equation}
C_{total} = C_{quant} + C_{trees} + C_{vote}
\label{eq:total-cycles}
\end{equation}

\noindent dimana:

\begin{itemize}
    \item $C_{quant}$ adalah siklus untuk kuantisasi fitur
    \item $C_{trees}$ adalah siklus untuk traversal semua pohon
    \item $C_{vote}$ adalah siklus untuk penghitungan voting
\end{itemize}

Siklus kuantisasi fitur dihitung sebagai:

\begin{equation}
C_{quant} = N_f \times C_{per\_feature} + C_{loop\_overhead}
\label{eq:quant-cycles}
\end{equation}

\noindent dengan $N_f$ adalah jumlah fitur dan $C_{per\_feature} \approx 35$ siklus (mencakup operasi muat skala, perkalian, pergeseran, penjumlahan, dan penyimpanan).

Siklus traversal pohon dihitung sebagai:

\begin{equation}
C_{trees} = N_t \times (C_{call} + L_{avg} \times C_{per\_node})
\label{eq:tree-cycles}
\end{equation}

\noindent dengan $N_t$ adalah jumlah pohon, $C_{call} \approx 10$ siklus untuk \textit{overhead} pemanggilan fungsi, $L_{avg} = \lfloor D_{max}/2 \rfloor + 1$ adalah panjang jalur rata-rata, dan $C_{per\_node} \approx 25$ siklus per node (mencakup operasi muat indeks fitur, muat threshold, perbandingan, percabangan, dan muat pointer anak).

\subsubsection{Keunggulan dan Keterbatasan}

Metode \textit{static cycle counting} memiliki beberapa keunggulan dibandingkan pengukuran dinamis:

\begin{enumerate}
    \item \textbf{Reprodusibilitas}: Hasil analisis bersifat deterministik dan dapat direproduksi tanpa variasi pengukuran.
    \item \textbf{Tidak Memerlukan Perangkat Keras}: Analisis dapat dilakukan pada sistem pengembangan tanpa memerlukan perangkat target fisik.
    \item \textbf{Analisis Worst-Case}: Memungkinkan estimasi \textit{Worst-Case Execution Time} (WCET) yang penting untuk sistem \textit{real-time}.
    \item \textbf{Granularitas Tinggi}: Dapat mengidentifikasi kontribusi siklus dari setiap bagian kode secara detail.
\end{enumerate}

Namun, metode ini juga memiliki keterbatasan:

\begin{enumerate}
    \item \textbf{Tidak Memperhitungkan Cache}: MSP430F5529 tidak memiliki cache instruksi, sehingga keterbatasan ini tidak relevan untuk platform target.
    \item \textbf{Estimasi Percabangan}: Untuk kode dengan banyak percabangan kondisional, diperlukan asumsi tentang distribusi jalur eksekusi.
    \item \textbf{Optimisasi Compiler}: Hasil bergantung pada level optimisasi compiler yang digunakan.
\end{enumerate}

\subsubsection{Konversi ke Waktu Eksekusi}

Setelah jumlah siklus diperoleh, waktu eksekusi dapat dihitung menggunakan frekuensi clock prosesor:

\begin{equation}
T_{inference} = \frac{C_{total}}{f_{clock}}
\label{eq:time-conversion}
\end{equation}

\noindent dimana $f_{clock} = 25$ MHz untuk MSP430F5529. Sebagai contoh, untuk model dengan $C_{total} = 1.270$ siklus:

\begin{equation}
T_{inference} = \frac{1.270}{25.000.000} = 50,8 \times 10^{-6} \text{ detik} = 50,8 \text{ }\mu\text{s}
\end{equation}

Throughput inferensi dapat dihitung sebagai:

\begin{equation}
\text{Throughput} = \frac{f_{clock}}{C_{total}} = \frac{25.000.000}{1.270} \approx 19.685 \text{ inferensi/detik}
\label{eq:throughput}
\end{equation}

Metode \textit{static cycle counting} ini memberikan estimasi yang akurat untuk evaluasi kelayakan deployment model \textit{machine learning} pada mikrokontroler dengan sumber daya terbatas, khususnya dalam konteks aplikasi WSN yang memiliki batasan waktu respons dan konsumsi daya yang ketat.
